% ============================================================
% eRUPT Academic Research Collaboration Proposal
% Template: eRUPT 2026/27
% ============================================================

\documentclass[11pt, letterpaper]{article}

% --- Page layout ---
\usepackage[letterpaper, margin=1in, top=0.75in, bottom=0.75in]{geometry}

% --- Font setup ---
\usepackage[utf8]{inputenc}
\usepackage[T1]{fontenc}
\usepackage{lmodern}
\usepackage{microtype}

% --- Packages ---
\usepackage{graphicx}
\usepackage{xcolor}
\usepackage{hyperref}
\usepackage{booktabs}
\usepackage{array}
\usepackage{tabularx}
\usepackage{amsmath, amssymb}
\usepackage{parskip}

% --- Colors from template ---
\definecolor{ebay-dark}{HTML}{353744}
\definecolor{ebay-green}{HTML}{00ab44}
\definecolor{ebay-gray}{HTML}{444444}

% --- Headings ---
\usepackage{titlesec}
\titleformat{\section}
  {\normalfont\bfseries\fontsize{14}{17}\selectfont\color{ebay-dark}}
  {\thesection}{0.5em}{}
\titleformat{\subsection}
  {\normalfont\bfseries\fontsize{12}{15}\selectfont\color{ebay-dark}}
  {\thesubsection}{0.5em}{}

% --- Title ---
\makeatletter
\renewcommand{\maketitle}{%
  {\parindent0pt\relax
   \fontsize{36}{43}\selectfont\color{ebay-dark}\bfseries\raggedright
   \@title\par}%
  {\color{gray!60}\hrule height 1.5pt\par}%
}
\makeatother

% --- Horizontal rule macro ---
\newcommand{\ehr}{{\color{gray!60}\hrule height 1.5pt}}

% --- Metadata ---
\title{Agentic Knowledge Graphs for\\Deterministic Code Generation}
\author{}
\date{}

% ============================================================
\begin{document}

\maketitle

% ------------------------------------------------------------
\section*{Engagement Structure}
\label{sec:engagement}
% ------------------------------------------------------------

\noindent
\begin{tabularx}{\textwidth}{>{\raggedright\arraybackslash}p{0.45\textwidth} X}
  \toprule
  \textbf{University Principal Investigator, Professor and Researcher(s)} &
  Dr.~Chris Brown (PI, Virginia Tech), Jason Cusati (GRA, Virginia Tech) \\
  \midrule
  \textbf{University} & Virginia Tech, USA \\
  \midrule
  \textbf{eBay Co-Investigator(s)} & Ramesh Periyathambi \\
  \midrule
  \textbf{eBay VP/Head of Department} & \\
  \midrule
  \textbf{Grant Amount Requested} & \\
  \bottomrule
\end{tabularx}

\ehr

% ------------------------------------------------------------
\section*{Objectives}
\label{sec:objectives}
% ------------------------------------------------------------

At eBay, adding a new signal to a ViewItem page or wiring up an A/B test variant is a task that takes developers 2--4 weeks, every time. The files are different, the APIs have shifted, the team conventions are undocumented---and someone has to rediscover all of it from scratch. Across a codebase of this scale, that pattern repeats hundreds of times a year, silently consuming engineering velocity.

The primary goal of this project is to make that pattern---and every recurring code pattern like it---deterministically automatable. We propose building an \textbf{agentic knowledge graph} (KG) that captures eBay's historical code patterns, team conventions, and dependency structures from git history, PRs, and code reviews. When a developer requests a new signal handler, an AI agent queries the KG for the most relevant prior implementations---exactly which files changed, which APIs were called, what review feedback came back---and generates code consistent with past practice. Same task, same KG state, same output.

Why academic research? The open challenges here---mining actionable patterns from a production codebase, representing them in a queryable KG, designing agentic workflows that use context deterministically, and building automated validation pipelines that check generated code against historical norms---require systematic investigation, not engineering alone. These span software engineering, AI, and knowledge representation in ways that no in-house team has the bandwidth to explore.

\ehr

% ------------------------------------------------------------
\section*{Previous Work}
\label{sec:previous}
% ------------------------------------------------------------

This proposal, developed in collaboration with Ramesh Periyathambi at eBay, builds on prior work by our team and by the broader research community.

\textbf{Compass (eBay Internal).} eBay's enterprise knowledge graph, Compass, already connects Jira tickets to git repos, commits, PRs, deployments, and microservices. Our agentic KG will extend this infrastructure with code-level pattern knowledge---what the tickets don't capture.

\textbf{PA-AKG (DOE Genesis).} Our ongoing DOE-funded project at Virginia Tech is building provenance-aware agentic knowledge graphs for scientific software [1]. The core architecture---AST-based extraction of code structure, agentic orchestration with ranking and validation agents, and provenance tracking---directly transfers to eBay's enterprise Java context.

\textbf{VersiCode.} The VersiCode dataset [2] demonstrated that version-aware code knowledge can be encoded in structured form, tracking API evolution across library releases---a capability directly relevant to eBay's ever-evolving dependency graph.

\textbf{FOSE 2026.} Our position paper at ICSE-FOSE 2026 [1] articulated the need for structured knowledge accumulation in software engineering---the intellectual foundation for this proposal.

\ehr

% ------------------------------------------------------------
\section*{Methodologies / Approaches}
\label{sec:method}
% ------------------------------------------------------------

We pursue two complementary strategies to mitigate risk:

\textbf{Approach 1: Code Knowledge Graph Construction \& Pattern Mining.} Working with Ramesh and his team at eBay, we will construct a knowledge graph from eBay's Java codebase using Abstract Syntax Trees (ASTs) to extract API usage, dependency relationships, and structural conventions. We augment this with historical data from git history and PR descriptions. Frequent sub-graphs across the codebase form the pattern library---for example, the ViewItem signal pattern, captured as a set of entity nodes (\texttt{ApiGateway}, \texttt{SignalHandler}, \texttt{MetricsClient}) with typed edges (\texttt{imports}, \texttt{implements}, \texttt{dependsOn}).

\textbf{Approach 2: Agentic Orchestration.} We deploy three role-specific agents:
\begin{itemize}
    \item \textit{Ranking agents} score candidate patterns from the KG by relevance to the current task, filtering noise.
    \item \textit{Generation agents} produce code conditioned on the ranked patterns, using the KG as structured context rather than relying solely on LLM priors.
    \item \textit{Validation agents} check generated PRs against KG-encoded patterns and flag deviations from team conventions or breaking dependency changes.
\end{itemize}

Together, these approaches make code generation deterministic---query the same KG state, get the same output---and auditable, because every recommendation traces back to historical patterns in the graph.

\ehr

% ------------------------------------------------------------
\section*{Impact}
\label{sec:impact}
% ------------------------------------------------------------

If successful, this project reduces the time for recurring code patterns from weeks to hours. The agentic KG becomes eBay's persistent organizational memory: new teams, offshore hires, and internal transfers can query it for project conventions, dependency ownership, and historical decisions without shadowing senior developers for weeks.

The validation pipeline directly addresses eBay's code-review bottleneck for AI-generated PRs. Rather than manually reviewing every agentic change, the KG provides automated conformance checking---flagging only genuine anomalies for human attention.

\textbf{Risks.} eBay's codebase is large; we mitigate by starting with well-scoped patterns (ViewItem signals) as case studies. KG extraction reliability is an open problem; we use staged evaluation with manual validation.

\ehr

% ------------------------------------------------------------
\section*{Expected Product / Output / Milestones}
\label{sec:output}
% ------------------------------------------------------------

\textbf{Deliverables:}
\begin{itemize}
    \item \textbf{PoC Agentic KG} (Months 1--7): Scoped to ViewItem signals, with pattern mining and agentic orchestration.
    \item \textbf{Validation Framework} (Months 8--10): Automated conformance checking for agent-generated PRs.
    \item \textbf{Empirical Evaluation} (Months 11--12): Developer time reduction, pattern retrieval accuracy, PR acceptance rates.
    \item \textbf{Publications:} Academic papers at ICSE, FSE, MSR disseminating findings.
\end{itemize}

\textbf{Metrics for success:} Pattern retrieval precision/recall; agent-generated code acceptance rate; reduction in developer time per pattern instance. Tiered funding could support progressive expansion to additional pattern types across eBay's organization.

\ehr

% ------------------------------------------------------------
\section*{University / Lab / Researcher Qualification}
\label{sec:qualification}
% ------------------------------------------------------------

\textbf{Dr.~Chris Brown (PI, Virginia Tech).} Dr.~Brown leads the Software Engineering and Digital Society lab, researching human-centered AI for software engineering and LLM-based code generation. He has published at ICSE, FSE, ESEM, and CSCW, and is PI on the DOE Genesis PA-AKG project whose architecture directly informs this proposal.

\textbf{Jason Cusati (GRA, Virginia Tech).} Mr.~Cusati is a PhD candidate whose research focuses on knowledge graphs for software engineering and AI-assisted code generation. He co-authored the FOSE 2026 paper on structured knowledge accumulation and is the lead researcher on the PA-AKG project.

\ehr

% ------------------------------------------------------------
\section*{References}
\label{sec:references}
% ------------------------------------------------------------

[1] J.~Cusati and C.~Brown. ``A Case for Structured Knowledge Accumulation in Software Engineering Research.'' In \textit{Proceedings of ICSE-FOSE}, 2026.

[2] T.~Wu et al. ``VersiCode: Towards Version-controllable Code Generation.'' arXiv:2406.07411, 2024.

\end{document}
